\documentclass[11pt,a4paper]{article}

% -------------------------------------------------
% Packages
% -------------------------------------------------
\usepackage[margin=2.2cm]{geometry}
\usepackage{kotex}
\usepackage{amsmath,amssymb}
\usepackage{booktabs}
\usepackage{array}
\usepackage{xcolor}
\usepackage{listings}
\usepackage[most]{tcolorbox}
\usepackage{enumitem}
\usepackage{hyperref}
\usepackage{fancyhdr}
\usepackage{titlesec}

% -------------------------------------------------
% Colors
% -------------------------------------------------
\definecolor{codebg}{RGB}{247,247,247}
\definecolor{codeframe}{RGB}{210,210,210}
\definecolor{keywordcolor}{RGB}{0,70,160}
\definecolor{commentcolor}{RGB}{0,120,80}
\definecolor{stringcolor}{RGB}{170,50,50}
\definecolor{titleblue}{RGB}{35,75,130}

% -------------------------------------------------
% Hyperlinks
% -------------------------------------------------
\hypersetup{
    colorlinks=true,
    linkcolor=titleblue,
    urlcolor=titleblue
}

% -------------------------------------------------
% Header / Footer
% -------------------------------------------------
\pagestyle{fancy}
\fancyhf{}
\lhead{Python Programming}
\rhead{Day 11}
\cfoot{\thepage}
\setlength{\headheight}{14pt}

% -------------------------------------------------
% Section style
% -------------------------------------------------
\titleformat{\section}
  {\Large\bfseries\color{titleblue}}
  {\thesection.}{0.6em}{}

\titleformat{\subsection}
  {\large\bfseries}
  {\thesubsection.}{0.5em}{}

% -------------------------------------------------
% Python code style
% -------------------------------------------------
\lstdefinestyle{pythonstyle}{
    language=Python,
    backgroundcolor=\color{codebg},
    frame=single,
    rulecolor=\color{codeframe},
    basicstyle=\ttfamily\small,
    keywordstyle=\color{keywordcolor}\bfseries,
    commentstyle=\color{commentcolor},
    stringstyle=\color{stringcolor},
    numbers=left,
    numberstyle=\tiny\color{gray},
    numbersep=8pt,
    tabsize=4,
    showstringspaces=false,
    breaklines=true,
    columns=fullflexible,
    keepspaces=true
}

\lstset{style=pythonstyle}

% -------------------------------------------------
% Boxes
% -------------------------------------------------
\newtcolorbox{learningbox}{
    colback=blue!4,
    colframe=titleblue,
    title=\textbf{학습 목표},
    fonttitle=\bfseries
}

\newtcolorbox{importantbox}{
    colback=yellow!8,
    colframe=orange!70!black,
    title=\textbf{핵심 포인트},
    fonttitle=\bfseries
}

\newtcolorbox{practicebox}{
    colback=red!3,
    colframe=red!55!black,
    title=\textbf{종합 실습},
    fonttitle=\bfseries
}

% -------------------------------------------------
% Document
% -------------------------------------------------
\begin{document}

\begin{titlepage}
    \centering
    \vspace*{3cm}

    {\Huge\bfseries Python Programming\par}
    \vspace{0.8cm}
    {\LARGE Day 11: Comprehensive Practice\par}

    \vspace{2cm}

    {\large Basics, Functions, Data Structures, OOP, NumPy, Pandas, Matplotlib\par}

    \vfill
\end{titlepage}

\tableofcontents
\newpage

% =================================================
\section{학습 목표}
% =================================================

\begin{learningbox}
Day 11은 Python Day 1--10에서 학습한 내용을 새로운 개념 없이 종합적으로 복습하는 과정이다.
이 장을 마치면 다음 내용을 직접 코드로 작성할 수 있다.
\begin{itemize}[leftmargin=1.5em]
    \item 조건문, 반복문, 리스트, 딕셔너리를 이용해 데이터를 처리하기
    \item 기능을 함수로 분리하고 매개변수와 반환값을 활용하기
    \item 중첩된 리스트와 딕셔너리를 탐색하고 수정하기
    \item 클래스, 객체, 상속, 오버라이딩을 이용해 데이터를 구조화하기
    \item NumPy 배열의 인덱싱, 슬라이싱, 배열 연산을 수행하기
    \item Pandas DataFrame을 생성하고 데이터를 선택, 수정, 분석하기
    \item Matplotlib을 이용해 데이터를 여러 형태의 그래프로 시각화하기
\end{itemize}
\end{learningbox}

\begin{importantbox}
Day 11의 목적은 새로운 문법을 추가하는 것이 아니라 이미 학습한 Python 기능을 직접 다시 구현하면서 기본기를 점검하는 것이다.
\end{importantbox}

% =================================================
\section{Basic Python}
% =================================================

첫 번째 실습에서는 학생 점수 데이터를 이용해 Python의 기본 자료구조와 제어문을 복습한다.
\texttt{sum()}, \texttt{max()}, \texttt{min()}, \texttt{sorted()}에 의존하지 않고 반복문과 조건문으로 직접 처리한다.

\subsection{최종 코드: \texttt{01\_basic.py}}

\begin{lstlisting}
students = [
    {"name": "Alice", "score": 85},
    {"name": "Bob", "score": 72},
    {"name": "Charlie", "score": 93},
    {"name": "David", "score": 68},
    {"name": "Emma", "score": 88}
]

total = 0

for i in range(0, len(students)):
    print(f'{students[i]["name"]}: {students[i]["score"]}')
    total += students[i]["score"]

print()

average = total / len(students)

print(f"Total: {total}")
print(f"Average: {average}")
print()

print("Students with score >= 80")

for i in range(0, len(students)):
    if students[i]["score"] >= 80:
        print(f'{students[i]["name"]}: {students[i]["score"]}')

print()

top_index = 0

for i in range(1, len(students)):
    if students[i]["score"] > students[top_index]["score"]:
        top_index = i

print(f'Top Student: {students[top_index]["name"]}')
print(f'Score: {students[top_index]["score"]}')
print()

for i in range(0, len(students)):
    if students[i]["score"] >= 90:
        grade = "A"
    elif students[i]["score"] >= 80:
        grade = "B"
    elif students[i]["score"] >= 70:
        grade = "C"
    elif students[i]["score"] >= 60:
        grade = "D"
    else:
        grade = "F"

    print(f'{students[i]["name"]}: {grade}')
\end{lstlisting}

\begin{importantbox}
최댓값을 직접 찾을 때는 첫 번째 원소를 현재 최고값으로 설정하고 두 번째 원소부터 비교하며 더 큰 값이 나오면 갱신할 수 있다.
\end{importantbox}

% =================================================
\section{Functions}
% =================================================

두 번째 실습에서는 기본 데이터 처리 로직을 함수로 분리한다.
함수 내부에서 모든 결과를 출력하기보다 필요한 값을 \texttt{return}으로 반환하여 호출한 코드에서 재사용한다.

\subsection{최종 코드: \texttt{02\_functions.py}}

\begin{lstlisting}
students = [
    {"name": "Alice", "score": 85},
    {"name": "Bob", "score": 72},
    {"name": "Charlie", "score": 93},
    {"name": "David", "score": 68},
    {"name": "Emma", "score": 88}
]


def calculate_total(students):
    total = 0

    for student in students:
        total += student["score"]

    return total


def calculate_average(students):
    return calculate_total(students) / len(students)


def get_high_scorers(students, cutoff):
    result = []

    for student in students:
        if student["score"] >= cutoff:
            result.append(student)

    return result


def get_top_student(students):
    top_index = 0

    for i in range(1, len(students)):
        if students[i]["score"] > students[top_index]["score"]:
            top_index = i

    return students[top_index]


def get_grade(score):
    if score >= 90:
        return "A"
    elif score >= 80:
        return "B"
    elif score >= 70:
        return "C"
    elif score >= 60:
        return "D"
    else:
        return "F"


print(f"Total: {calculate_total(students)}")
print(f"Average: {calculate_average(students)}")
print()

print("Students with score >= 80")

for student in get_high_scorers(students, 80):
    print(f'{student["name"]}: {student["score"]}')

print()

top_student = get_top_student(students)

print(f'Top Student: {top_student["name"]}')
print(f'Score: {top_student["score"]}')
print()

print("Grades")

for student in students:
    print(f'{student["name"]}: {get_grade(student["score"])}')
\end{lstlisting}

\begin{importantbox}
함수의 결과를 여러 번 사용할 경우 같은 함수를 반복 호출하기보다 반환값을 변수에 저장하여 재사용할 수 있다.
\end{importantbox}

% =================================================
\section{Data Structures}
% =================================================

세 번째 실습에서는 리스트 안에 딕셔너리가 있고, 그 딕셔너리 안에 다시 리스트와 딕셔너리가 들어 있는 중첩 자료구조를 다룬다.
학생 중심 데이터를 탐색하고 수정한 뒤 과목 중심 데이터로 다시 구성한다.

\subsection{최종 코드: \texttt{03\_data\_structures.py}}

\begin{lstlisting}
students = [
    {
        "name": "Alice",
        "major": "Physics",
        "courses": ["Python", "Math"],
        "scores": {"Python": 88, "Math": 92}
    },
    {
        "name": "Bob",
        "major": "Engineering",
        "courses": ["Python", "C++"],
        "scores": {"Python": 76, "C++": 85}
    },
    {
        "name": "Charlie",
        "major": "Physics",
        "courses": ["Math", "C++"],
        "scores": {"Math": 95, "C++": 91}
    },
    {
        "name": "David",
        "major": "Engineering",
        "courses": ["Python", "Math", "C++"],
        "scores": {"Python": 82, "Math": 79, "C++": 88}
    }
]

print("Physics Students")

for student in students:
    if student["major"] == "Physics":
        print(f'{student["name"]}: {student["courses"]}')

print()

print("Python Students")

for student in students:
    for course in student["courses"]:
        if course == "Python":
            print(f'{student["name"]}: {student["scores"]["Python"]}')

print()

for student in students:
    if student["name"] == "Bob":
        student["courses"].append("Math")
        student["scores"]["Math"] = 84
        print(student)

print()

new_student = {
    "name": "Emma",
    "major": "Physics",
    "courses": ["Python", "C++", "Math"],
    "scores": {"Python": 94, "C++": 90, "Math": 89}
}

students.append(new_student)

for student in students:
    print(student["name"])

print()

cpp_list = []

for student in students:
    for course in student["courses"]:
        if course == "C++":
            cpp_list.append(
                {
                    "name": student["name"],
                    "score": student["scores"]["C++"]
                }
            )

top_index = 0

for i in range(1, len(cpp_list)):
    if cpp_list[i]["score"] > cpp_list[top_index]["score"]:
        top_index = i

print(f'Top C++ Student: {cpp_list[top_index]["name"]}')
print(f'Score: {cpp_list[top_index]["score"]}')
print()

course_students = {}
py_group = []
cpp_group = []
math_group = []

for student in students:
    for course in student["courses"]:
        if course == "Python":
            py_group.append(student["name"])
        elif course == "C++":
            cpp_group.append(student["name"])
        elif course == "Math":
            math_group.append(student["name"])

course_students["Python"] = py_group
course_students["Math"] = math_group
course_students["C++"] = cpp_group

print(course_students)
\end{lstlisting}

\begin{importantbox}
중첩 자료구조에서는 현재 접근하고 있는 값이 리스트인지 딕셔너리인지 구분하는 것이 중요하다.
리스트에는 인덱스와 \texttt{append()}를 사용하고, 딕셔너리에는 key를 이용해 값을 읽거나 추가할 수 있다.
\end{importantbox}

% =================================================
\section{Object-Oriented Programming}
% =================================================

네 번째 실습에서는 영화 데이터를 클래스로 표현한다.
기본 클래스와 자식 클래스를 만들고, 객체들을 하나의 리스트에 저장하여 상속과 오버라이딩을 복습한다.

\subsection{최종 코드: \texttt{04\_oop.py}}

\begin{lstlisting}
class Movie:
    def __init__(self, title, year, rating):
        self.title = title
        self.year = year
        self.rating = rating

    def show_info(self):
        print(f"Title: {self.title}")
        print(f"Year: {self.year}")
        print(f"Rating: {self.rating}")

    def is_high_rating(self):
        return self.rating >= 9.0


movies = []

movies.append(Movie("Interstellar", 2014, 9.1))
movies.append(Movie("Inception", 2010, 8.8))
movies.append(Movie("Parasite", 2019, 8.6))
movies.append(Movie("The Dark Knight", 2008, 9.0))

for movie in movies:
    movie.show_info()
    print()

print("High Rated Movies")

for movie in movies:
    if movie.is_high_rating():
        print(movie.title)

print()


class StreamingMovie(Movie):
    def __init__(self, title, year, rating, platform):
        super().__init__(title, year, rating)
        self.platform = platform

    def show_info(self):
        super().show_info()
        print(f"Platform: {self.platform}")


movies.append(StreamingMovie("Dune", 2021, 8.5, "Netflix"))
movies.append(StreamingMovie("Oppenheimer", 2023, 8.9, "Watcha"))

for movie in movies:
    movie.show_info()
    print()

top_rating_movie = movies[0]

for i in range(1, len(movies)):
    if movies[i].rating > top_rating_movie.rating:
        top_rating_movie = movies[i]

print("Top Movie")
print(f"Title: {top_rating_movie.title}")
print(f"Rating: {top_rating_movie.rating}")
\end{lstlisting}

\begin{importantbox}
부모 클래스의 생성자와 메서드는 \texttt{super()}를 이용해 재사용할 수 있다.
같은 \texttt{show\_info()} 호출이라도 실제 객체가 자식 클래스의 객체라면 오버라이딩된 메서드가 실행된다.
\end{importantbox}

% =================================================
\section{NumPy}
% =================================================

다섯 번째 실습에서는 NumPy 배열 생성, 인덱싱, 슬라이싱, 원소별 배열 연산과 2차원 배열 접근을 복습한다.

\subsection{최종 코드: \texttt{05\_numpy.py}}

\begin{lstlisting}
import numpy as np


a = np.array([10, 20, 30, 40, 50])
b = np.array([2, 4, 6, 8, 10])

matrix = np.array([
    [1, 2, 3],
    [4, 5, 6],
    [7, 8, 9]
])

print(a[0])
print(a[-1])
print(a[2])
print()

print(a[1:4])
print()

print(a + b)
print(a - b)
print(a * b)
print(a / b)
print()

print(a + 5)
print()

print(b * 3)
print()

print(matrix[1, 1])
print()

print(matrix[1, :])
print()

print(matrix * 2)
print(matrix + 10)
print()

print(np.zeros(5))
print(np.ones(5))
print(np.arange(0, 10, 2))
print(np.linspace(0, 1, 5))
\end{lstlisting}

\begin{importantbox}
NumPy 배열에 대한 \texttt{+}, \texttt{-}, \texttt{*}, \texttt{/} 연산은 같은 위치의 원소에 대해 수행된다.
스칼라와 배열을 연산하면 해당 연산이 배열의 모든 원소에 적용된다.
\end{importantbox}

% =================================================
\section{Pandas}
% =================================================

여섯 번째 실습에서는 딕셔너리로부터 DataFrame을 생성하고 열과 행을 선택한 뒤 새로운 열을 추가하고 기본 통계를 확인한다.

\subsection{최종 코드: \texttt{06\_pandas.py}}

\begin{lstlisting}
import pandas as pd


data = {
    "Name": ["Alice", "Bob", "Charlie", "David", "Emma"],
    "Physics": [88, 76, 95, 82, 91],
    "Math": [92, 85, 90, 78, 94]
}

df = pd.DataFrame(data)

print(df)
print()

print(df["Physics"])
print(df["Name"])
print()

print(df.loc[2])
print(df.loc[1:3, ["Name", "Physics"]])
print()

print(df.iloc[0])
print(df.iloc[-1])
print()

df["Average"] = (df["Physics"] + df["Math"]) / 2

print(df)
print()

print(f"Average of Physics: {df['Physics'].mean()}")
print(f"Average of Math: {df['Math'].mean()}")
print(f"Maximum of Physics: {df['Physics'].max()}")
print(f"Minimum of Math: {df['Math'].min()}")
print()

print(df.describe())
\end{lstlisting}

\begin{importantbox}
\texttt{loc}는 label을 기준으로 데이터를 선택하고 \texttt{iloc}는 정수 위치를 기준으로 선택한다.
DataFrame의 열끼리 직접 연산하여 새로운 열을 만들 수도 있다.
\end{importantbox}

% =================================================
\section{Matplotlib}
% =================================================

일곱 번째 실습에서는 NumPy 배열을 Matplotlib에 전달하여 선 그래프, 산점도, 막대그래프를 그린다.
제목, 축 이름, 범례, 격자를 추가하여 기본적인 데이터 시각화 흐름을 복습한다.

\subsection{최종 코드: \texttt{07\_matplotlib.py}}

\begin{lstlisting}
import numpy as np
import matplotlib.pyplot as plt


x = np.array([1, 2, 3, 4, 5])

physics = np.array([70, 75, 83, 88, 95])
math = np.array([65, 72, 80, 85, 90])

plt.plot(x, physics)
plt.title("Physics Scores")
plt.xlabel("Test")
plt.ylabel("Score")
plt.show()

plt.plot(x, physics, label="Physics")
plt.plot(x, math, label="Math")
plt.title("Physics vs Math")
plt.xlabel("Test")
plt.ylabel("Score")
plt.grid(True)
plt.legend()
plt.show()

plt.scatter(physics, math)
plt.title("Physics vs Math Scores")
plt.xlabel("Physics Scores")
plt.ylabel("Math Scores")
plt.show()

students = ["Alice", "Bob", "Charlie", "David", "Emma"]

plt.bar(students, physics)
plt.title("Physics Scores by Student")
plt.xlabel("Student")
plt.ylabel("Physics Score")
plt.show()

x = np.linspace(0, 10, 100)
y = x ** 2

plt.plot(x, y)
plt.title("y = x^2")
plt.xlabel("x")
plt.ylabel("y")
plt.grid(True)
plt.show()
\end{lstlisting}

\begin{importantbox}
NumPy는 수치 데이터를 생성하고 계산하는 역할을 담당하며 Matplotlib은 계산된 데이터를 그래프로 표현할 수 있다.
\texttt{plt.show()}를 각 그래프 뒤에 호출하면 그래프를 하나씩 분리하여 확인할 수 있다.
\end{importantbox}

% =================================================
\section{종합 실습 파일 목록}
% =================================================

\begin{practicebox}
Day 11의 실습 파일은 다음과 같이 관리한다.
\begin{itemize}[leftmargin=1.5em]
    \item \texttt{01\_basic.py}
    \item \texttt{02\_functions.py}
    \item \texttt{03\_data\_structures.py}
    \item \texttt{04\_oop.py}
    \item \texttt{05\_numpy.py}
    \item \texttt{06\_pandas.py}
    \item \texttt{07\_matplotlib.py}
    \item \texttt{python\_day11.tex}
    \item \texttt{python\_day11.pdf}
\end{itemize}
\end{practicebox}

% =================================================
\section{Day 11 핵심 요약}
% =================================================

\begin{itemize}[leftmargin=1.5em]
    \item 리스트와 딕셔너리는 여러 데이터를 구조적으로 저장하고 반복문으로 처리할 수 있다.
    \item 조건문을 이용하면 데이터의 값에 따라 서로 다른 처리를 수행할 수 있다.
    \item 최댓값을 직접 찾을 때는 하나의 값을 현재 최고값으로 두고 순회하면서 갱신할 수 있다.
    \item 함수는 반복되는 기능을 분리하고 매개변수를 통해 데이터를 전달받으며 \texttt{return}으로 결과를 반환한다.
    \item 함수의 반환값은 변수에 저장하여 여러 번 재사용할 수 있다.
    \item 중첩 자료구조에서는 각 단계의 자료형을 구분하며 접근해야 한다.
    \item 클래스는 데이터와 관련 기능을 하나의 객체로 묶을 수 있다.
    \item \texttt{\_\_init\_\_()}은 객체가 생성될 때 인스턴스 변수를 초기화한다.
    \item 상속을 이용하면 기존 클래스의 기능을 재사용하고 새로운 기능을 추가할 수 있다.
    \item 메서드 오버라이딩을 이용하면 자식 클래스에서 부모 클래스의 동작을 변경하거나 확장할 수 있다.
    \item NumPy 배열은 인덱싱과 슬라이싱을 지원하며 배열 전체에 효율적으로 수치 연산을 적용할 수 있다.
    \item 2차원 NumPy 배열은 행과 열의 인덱스를 이용해 원하는 원소와 부분 배열에 접근할 수 있다.
    \item Pandas DataFrame은 행과 열 형태의 데이터를 다루기 위한 자료구조이다.
    \item \texttt{loc}와 \texttt{iloc}를 이용하면 원하는 행과 열을 선택할 수 있다.
    \item DataFrame의 열끼리 연산하여 새로운 열을 만들고 \texttt{mean()}, \texttt{max()}, \texttt{min()}, \texttt{describe()}로 데이터를 분석할 수 있다.
    \item Matplotlib의 \texttt{plot()}, \texttt{scatter()}, \texttt{bar()}를 이용해 여러 형태의 그래프를 만들 수 있다.
    \item NumPy로 생성하거나 계산한 데이터를 Matplotlib으로 시각화할 수 있다.
    \item Day 11은 Python Day 1--10에서 학습한 내용을 영역별 종합 실습으로 다시 확인하는 과정이다.
\end{itemize}

\end{document}
